In this section we record the combinatorial and algebraic structure of the doubly pruned simplicial complexes themselves. We show that every such complex is vertex decomposable, hence shellable and sequentially Cohen-Macaulay. We also describe its facets explicitly, which gives a criterion for when $\Delta_w$ fails to be pure. We use this description in \Cref{subsec:hilb} to compute the Hilbert series of both $R_{\Delta_w}$ and $R_{\Delta_w,n}.$
 
For the rest of this section, write $\supp(\Delta) = \{v: \{v\}\in\Delta\}$ for the set of vertices of $\Delta.$ 
Recall the two subcomplexes that play a role in \Cref{def: vertex-decomposable}, that is $\lk_v(\Delta)$ and $\del_v(\Delta),$ the link and deletion of $v$ in $\Delta,$ respectively. 
Since $\lk_v(\Delta)\subseteq\del_v(\Delta),$ the shedding condition takes a simpler form.

\begin{remark}
    A vertex $v$ is a shedding vertex of $\Delta$ if and only if $G\cup\{v\}\not\in\Delta$ for every facet $G$ of $\del_v(\Delta).$ In particular, if $G$ is such a facet with $G\cup\{v\}\in\Delta,$ then $G\in\lk_v(\Delta),$ and any facet of $\lk_v(\Delta)$ containing $G$ also lies in $\del_v(\Delta)$ and hence equals $G$, so $G$ is a facet of both. The converse is immediate.  
\end{remark}


\begin{lemma}
\label{lem: two-min faces implies vertex-decomposable}
Let $\Delta$ be a simplicial complex on a finite set $V$ with at most two minimal nonfaces of cardinality at least two. Then $\Delta$ is vertex decomposable.
\end{lemma}
\begin{proof}
    We proceed by induction on the cardinality of the vertex set $V$ of $\Delta,$ with $\Delta = \varnothing$ or $\Delta$ being a simplex (including the singletons) being clear.
    If some minimal nonface is a singleton $\{u\},$ then no face of $\Delta$ contains $u,$ so $\Delta$ is also a simplicial complex on the smaller vertex set $V\setminus\{u\}$ with the same faces, and whether $\Delta$ is vertex decomposable does not depend on which of these vertex sets it is regarded. Thus, we may discard $u$ and so assume that no minimal nonface is a singleton.
    
    Now assume every minimal nonface has at least two elements. 
    Call them $A$ and $B$ (if a second minimal nonface exists). 
    If there are two minimal nonfaces with $A\cap B\neq\varnothing,$ choose $v\in A\cap B;$ otherwise choose $v\in A.$ 
    It follows that $v$ then lies in every minimal nonface unless $A$ and $B$ exist and are disjoint.
    If $v$ lives in every minimal nonface, then no subset of $V\setminus\{v\}$ contains a minimal nonface, so $\del_v(\Delta)$ is the simplex on $V\setminus\{v\},$ whose unique facet is $G = V\setminus\{v\}$ and $G\cup\{v\} = V\supseteq A$ is a nonface. 
    Otherwise, $A\cap B = \varnothing$ and $v\in A.$ 
    So, $\del_v(\Delta) = \{F\subseteq V\setminus\{v\}:B\not\subseteq F\},$ whose facets are the sets $G = V\setminus\{v,b\}$ with $b\in B.$  
    Since $b\not\in A,$ then $A\subseteq V\setminus\{b\} = G\cup\{v\}$ is a nonface. 
    In the first case, $\del_v(\Delta)$ is a simplex and has no minimal nonfaces at all, while in the second case $\del_v(\Delta) = \{F\subseteq V\setminus\{v\}:B\not\subseteq F\}$ has $B$ as its unique minimal nonface. So in both cases, $\del_v(\Delta)$ has at most one minimal nonface, and it is vertex decomposable by induction. Moreover, in both cases we have shown that $G\cup\{v\}\not\in\Delta$ for every facet $G$ of $\del_v(\Delta),$ so $v$ is a shedding vertex.
    Finally, $F\in\lk_v(\Delta)$ if and only if $F\cup\{v\}\in\Delta,$ i.e., if and only if $N\setminus\{v\}\not\subseteq F$ for every minimal nonface $N$ of $\Delta.$ So the minimal nonfaces of $\lk_v(\Delta)$ are the minimal elements of $\{A\setminus\{v\}, B\setminus\{v\}\}.$ Again, there are at most two, so $\lk_v(\Delta)$ is vertex-decomposable by induction. 
\end{proof}
\begin{theorem}
\label{thm: doubly pruned complexes imply vertex decomposable}
    Every doubly pruned simplicial complex is vertex decomposable. In particular, for every $w\in\sn(\calP),$ the complex $\Delta_w$ is vertex-decomposable, implying they are both shellable and sequentially Cohen-Macaulay.
\end{theorem}
\begin{proof}
    The minimal nonfaces of a doubly pruned complex of cardinality at least two are among $[i,p]$ and $[q,j],$ so \Cref{lem: two-min faces implies vertex-decomposable} applies. The remaining implications are immediate. 
\end{proof}
It is also possible to describe the facets of these complexes, which is what we do in \Cref{prop: facet description of doubly pruned complexes}. For a doubly pruned complex, normalize, so that $i<p$ and $q<j,$ which gives us the interval $[i,j].$ For $\Delta_w$ with $w\neq\id$ it is $[i(w),j(w)],$ since $p(w) = i(w)$ or $q(w) = j(w)$ would make $i(w)$ or $j(w)$ a fixed point of $w.$
\begin{proposition}
\label{prop: facet description of doubly pruned complexes}
Let $\Delta$ be a simplicial complex on vertex set $V$ with $W = \supp(\Delta)$ and with at most two minimal nonfaces $A,B$ of cardinality of at least two. 
Then $F$ is a facet of $\Delta$ if and only if $W\setminus F$ is a minimal subset of $W$ meeting each of $A$ and $B.$ Explicitly, when both $A$ and $B$ exist, the facets of $\Delta$ are 
\[
W\setminus\{v\}\quad (v\in A\cap B)\quad\text{and}\quad W\setminus\{a,b\}\quad (a\in A\setminus B,b\in B\setminus A).
\]
If only $A$ exists, the facets are $W\setminus\{v\}$ for $v\in A,$ and if neither exists, the unique facet is $W.$
Consequently, $\Delta$ is pure unless both $A$ and $B$ exist and $A\cap B\neq\varnothing,$ in which case $\Delta$ has facets of dimension $|W|-2$ and of dimension $|W|-3$ and $\dim(\Delta) = |W|-2.$
\end{proposition}
\begin{proof}
    Every minimal nonface is contained in $W$ or is a singleton outside of $W.$ 
    So for $F\subseteq W,$ by setting $T = W\setminus F,$ we have $F\in\Delta$ if and only if $T$ meets each of $A$ and $B.$ Complementation reverses the inclusion, which gives the first claim.
    A singleton $\{v\}$ meets both exactly when $v\in A\cap B,$ and singletons are automatically minimal. If $T$ is minimal with $|T|\geq 2,$ then $T\cap A\cap B = \varnothing,$ since one element of $T\cap A\cap B$ would already suffice. By choosing $a\in T\cap A$ and $b\in T\cap B,$ we obtain $a\in A\setminus B,b\in B\setminus A$ with $a\neq b,$ and $\{a,b\}\subseteq T$ meets both, so $T = \{a,b\}.$  Conversely, each such pair is minimal, as $\{a\}$ misses $B$ and \{b\} misses $A.$ 
    If only $A$ exists, then $T$ meets $A$ and $B$ exactly when $T\cap A\neq\varnothing,$ whose minimal solutions are the singletons $\{v\},$ $v\in A;$ if neither exists, $T = \varnothing$ is the unique minimal solution. Finally, $A$ and $B$ are incomparable, so whenever both exist, a facet of the second kind exists while a facet of the first kind exists exactly when $A\cap B\neq\varnothing.$
\end{proof}
\begin{corollary}
\label{cor: purity of doubly pruned complexes}
    Let $w\in\sn(\calP)$ with $w\neq\id$ and put $m = j(w) - i(w) + 1.$ Then $\Delta_w$ is not pure precisely when $i(w)<q(w)\leq p(w)<j(w).$ In this case, $\dim\Delta_w = m - 2$ and $\Delta_w$ has facets of dimension $m - 2$ and $m - 3.$
    If $q(w) = i(w)$ or $p(w) = j(w),$ then $\Delta_w$ is pure of dimension $m - 2.$ Otherwise, when $p(w) < q(w),$ it is pure of dimension $m - 3.$
\end{corollary}
\begin{proof}
    Here $W = [i(w),j(w)]$ and the minimal nonfaces of cardinality at least two are the minimal members of $\{[i(w),p(w)],[q(w),j(w)]\}.$ One of these intervals contains the other exactly when $q(w) = i(w)$ or $p(w) = j(w),$ leaving a single minimal nonface. Otherwise, both are minimal and meet in the interval $[q(w),p(w)]$ when $q(w)<p(w)$.
    To finish we apply \Cref{prop: facet description of doubly pruned complexes}.
\end{proof}

\begin{remark}
    \label{rem: doubly pruned purity}
    If a doubly pruned complex $\Delta$ is not $\Delta_w$ for some $w\in S_n(\calP)$, then it is pure if and only if $\Delta = \mathscr{P}([i,j])$, in which case it has dimension $j-i$. 
    If it is not pure, it has dimension $j-i-1$.
    In the notation of the corollary, these have dimensions $m-1$ and $m-2$, respectively.
\end{remark}


\begin{example}
\label{ex: 3412}
    Let $w = 3412 \in S_4(\calP)$, so that $w^{-1} = 3412$ as well. Its non-fixed
    points are $1,2,3,4$, so $i(w) = 1$ and $j(w) = 4$, giving $m = 4$, and
    \[
    p(w) = w^{-1}(i(w)) = w^{-1}(1) = 3,
    \qquad
    q(w) = w^{-1}(j(w)) = w^{-1}(4) = 2.
    \]
    By \Cref{def:Deltaw} the minimal nonfaces of $\Delta_w$ are the two intervals
    \[
    A = [i(w),p(w)] = \{1,2,3\}
    \qquad\text{and}\qquad
    B = [q(w),j(w)] = \{2,3,4\},
    \]
    with no singletons among them since $[i(w),j(w)] = [4]$. Here
    $A\cap B = \{2,3\}$, $A\setminus B = \{1\}$ and $B\setminus A = \{4\}$, so by
    \Cref{prop: facet description of doubly pruned complexes} the minimal subsets of
    $[4]$ meeting both $A$ and $B$ are $\{2\}$, $\{3\}$ and $\{1,4\}$, and the facets
    of $\Delta_w$ are their complements
    \[
    \{1,3,4\},\qquad \{1,2,4\},\qquad \{2,3\},
    \]
    of dimensions $2,2$ and $1$. Thus $\Delta_w$ is not pure and
    $\dim\Delta_w = 2$, in agreement with
    \Cref{cor: purity of doubly pruned complexes}, since
    $i(w) = 1 < 2 = q(w) \leq p(w) = 3 < 4 = j(w)$ and $m-2 = 2$. Pictorially,
     \begin{center}
    \begin{tikzpicture}[scale=0.4, vtx/.style={circle,fill=black,inner sep=1.7pt}]
      \coordinate (v1) at (0,2.6);
      \coordinate (v3) at (-2.2,-1);
      \coordinate (v2) at (2.2,-1);
      \coordinate (v4) at (0,0.45);
      \fill[faceblue] (v1) -- (v3) -- (v4) -- cycle;
      \fill[faceblue] (v1) -- (v2) -- (v4) -- cycle;
      \draw[thick] (v1) -- (v3) -- (v2) -- cycle;
      \draw[thick] (v1) -- (v4);
      \draw[thick] (v3) -- (v4);
      \draw[thick] (v2) -- (v4);
      \node[vtx] at (v1) {}; \node[vtx] at (v2) {};
      \node[vtx] at (v3) {}; \node[vtx] at (v4) {};
      \node[above]                     at (v1) {\small $1$};
      \node[below right]               at (v2) {\small $2$};
      \node[below left]                at (v3) {\small $3$};
      \node[above right,inner sep=2pt] at (v4) {\small $4$};
      \node[anchor=east] at (-2.75,0.8) {\large $\Delta_w \;=$};
    \end{tikzpicture}
    \end{center}
    Following the proof of
    \Cref{lem: two-min faces implies vertex-decomposable}, the shedding vertex is any $v \in A\cap B$, say $v = 2$. Then $\del_2(\Delta_w)$ is the simplex on $\{1,3,4\}$, while $\lk_2(\Delta_w)$ has minimal nonfaces $A\setminus\{2\} = \{1,3\}$ and $B\setminus\{2\} = \{3,4\}$ and hence facets $\{1,4\}$ and $\{3\}$. The unique facet $\{1,3,4\}$ of $\del_2(\Delta_w)$ satisfies $\{1,3,4\}\cup\{2\} = [4]\supseteq A$, a nonface, confirming that $2$ is a shedding vertex.
\end{example}
\begin{comment}
Alternate.
\begin{lemma}
\label{lem:pure-intersection}
    Let $u,v\in\sn(\calP).$ If $\Xi = \Delta_u\cap\Delta_v$ is pure of dimension $\dim\Delta_u,$ then $\Delta_u$ is pure.
\end{lemma}
\begin{proof}
    The nonfaces of $\Xi$ are generated by the union of the two generating sets of doubly pruned complexes. Set $I = \max\{i(u),i(v)\}$ and $J = \min\{j(u),j(v)\}.$ The singleton generators are then the $\{k\}$ with $k\not\in[I,J].$ Note that the four interval generators $[i(u),p(u)],[i(v),p(v)],[q(u),j(u)],[q(v),j(v)]$ have left endpoint $i(u)$ or $i(v),$ (respectively right endpoint $j(u)$ or $j(v)).$ Any generator with left endpoint smaller than $I$ or right endpoint larger than $J$ contains one of the singleton generators and may be discarded. Each surviving generator, therefore, lies in the interval $[I,J]$ and has left endpoint $I$ or right endpoint $J.$ Of the two surviving generators with left endpoint $I,$ one contains the other. It is analogous for the surviving two generators with right endpoint $J.$ What remains are the singletons outside of $[I,J]$ together with at most one interval $[I,p]$ and at most one interval $[q,J].$ Thus, $\Xi$ is again doubly pruned. Now set $W = \supp(\Delta_u) = [i,j], W' = \supp(\Xi = [i',j'], m = |W|,$ and $m' = |W|'.$ Since $\Xi\subseteq\Delta_u,$ then $W'\subseteq W$ and $m'\leq m.$
    Suppose $\Delta_u$ is not pure. By \Cref{prop: facet description of doubly pruned complexes}, $\Delta_u$ has two nonsingleton minimal nonfaces $A = [i,p]$ and $B = [q,j]$ with $A\cap B\neq\varnothing,$ and $\dim\Delta_u = m - 2.$ 
    Minimality forces $i<q$ and $p<j.$ Again, by \Cref{prop: facet description of doubly pruned complexes}, the pure complex $\Xi$ has dimension $m'-1,m'-2,$ or $m'-3,$ since it has zero, one, or two minimal nonfaces of cardinality at least two. 
    Since $m' - 3 = m - 2,$ would give $m' = m + 1 > m,$ then this case is impossible, so only two cases remain. 
    If $\Xi$ has one minimal nonface $C$ and $m' = m,$ then $C\subseteq A$ and $C\subseteq B,$ so $C\subseteq [q,p].$ However, $C$ has left endpoint $i$ or right endpoint $j,$ so $i\in[q,p]$ or $j\in[q,p].$ This forces $q\leq i$ or $j\leq p$ and contradicts $i<q$ and $p<j.$ If $\Xi$ is a simplex on $W'$ and $m' = m - 1,$ then $W' = [i+1,j]$ or $W' = [i,j-1],$ and the minimal nonfaces of $\Xi$ are the singletons outside $W'.$ But both $A$ and $B$ must contain one of them, which forces $q = i$ in the first case and $p = j$ in the second, once again forcing a contradiction.   
\end{proof}
\end{comment}

\begin{lemma}
\label{lem:pure-intersection}
    Let $u,v\in\sn(\calP).$ If $\Omega = \Delta_u\cap\Delta_v$ is pure of dimension $\dim\Delta_u,$ then $\Delta_u$ is pure.
\end{lemma}
\begin{proof}
    By \Cref{prop:intersection is doubly pruned}, $\Omega$ is doubly pruned. Set $W = \supp(\Delta_u) = [i,j]$, $W' = \supp(\Omega) = [i',j']$, $m = |W|$, and $m' = |W'|.$ Since $\Omega\subseteq\Delta_u,$ then $W'\subseteq W$ and $m'\leq m.$

    Suppose $\Delta_u$ is not pure, and write $A = [i,p]$ and $B = [q,j]$ for its two minimal nonfaces of cardinality at least two. Then $u\neq\id$, so \Cref{cor: purity of doubly pruned complexes} gives $i<q\leq p<j$ and $\dim\Delta_u = m-2.$ By \Cref{prop: facet description of doubly pruned complexes}, the pure complex $\Omega$ has dimension $m'-1,m'-2,$ or $m'-3,$ since it has zero, one, or two minimal nonfaces of cardinality at least two. Since $m' - 3 = m - 2$ would give $m' = m + 1 > m,$ this case is impossible, so only two cases remain.

    If $\dim(\Omega) = m'-2$, then it has one minimal nonface $C$ of cardinality at least 2 and $m' = m,$ so $C\subseteq A$ and $C\subseteq B,$ which implies that $C\subseteq [q,p].$ However, $C$ has left endpoint $i$ or right endpoint $j,$ so $i\in[q,p]$ or $j\in[q,p].$ This forces $q\leq i$ or $j\leq p$ and contradicts $i<q$ and $p<j.$ If $\Omega$ is a simplex on $W'$ and $m' = m - 1,$ then $W' = [i+1,j]$ or $W' = [i,j-1],$ and the minimal nonfaces of $\Omega$ are the singletons outside $W'.$ But both $A$ and $B$ must contain one of them, which forces $q = i$ in the first case and $p = j$ in the second, once again forcing a contradiction.
\end{proof}


%------------------------------------------------------------
\subsection{The Hilbert series for the Stanley-Reisner rings of doubly pruned simplicial complexes}
\label{subsec:hilb}
%------------------------------------------------------------

Let $S = \bK[x_1,\ldots,x_n]$, let $w\in\sn(\calP)$, and let $I_{\Delta_w}$ be the Stanley-Reisner ideal of $\Delta_w.$ 
In this section, we give formulations of the Hilbert series for both $R_{\Delta_w} = S/I_{\Delta_w}$ and of $R_{\Delta_w,n} = S/(I_{\Delta_w}+I_n),$ where $I_n$ is generated by the elementary symmetric functions.
Recall the definition of $\Delta_w$ for $w\in S_n(\calP)$ from Definition~\ref{def:Deltaw}. As in earlier sections, we write $i = i(w), j = j(w), p = p(w),$ and $q = q(w).$ We also set $A := [i,p]$ and $B := [q,j],$ two interval nonfaces of $\Delta_w.$
When $A$ and $B$ meet, we write $C := A\cap B$ and $\gamma = := \#A + \#B - m,$ so that $\#C = \#A\cap B.$
The Stanley-Reisner ideal is then
\[
I_{\Delta_w} = (x_1,\ldots,x_{i-1},x_{j+1},\ldots,x_n,x_{A},x_{B}),
\]
where $x_{[c,d]} = x_cx_{c+1}\cdots x_d.$
We now define three parameters used throughout this section:
\begin{itemize}
    \item $m = j - i + 1$, the size of the support of $\Delta$,
    \item $\#A = p - i + 1 = \#[i,p],$
    \item $\#B = j - q + 1 = \#[q,j].$
\end{itemize}
The two interval nonfaces are disjoint precisely when $p < q$ (equivalently $\#A + \#B \leq m$) they overlap when $q \leq p$ (equivalently $\#A + \#B > m$), in which case their overlap has size $\#A + \#B - m.$ 
If neither interval contains the other, then in the overlapping case we have $A\cup B = [i,j],$ i.e. the union already used the realized vertices. One interval contains the other exactly when $q = i$ or $p = j,$ equivalently when $\#B = m$ or $\#A = m;$ by \Cref{cor: purity of doubly pruned complexes} this is exactly the case in which $\Delta_w$ has a single minimal nonface of cardinality at least two, and we call it the \emph{nested case.} In the nested case, we write $[c,d]$ for that nonface, so $[c,d] = A$ if $\#B = m$ and $[c,d] = B$ if $\#A = m,$ and we set $k = \#[c,d] = \min(\#A,\#B).$ The nested case is not exotic as it contains every simple transposition $s_r,$ for which $[i,j] = [r,r+1]$ and $p = j,$ and it is the majority of $\sn(\calP)$ for small $n.$ It is, however, the easy case, since then $R_{\Delta_w} = S'/(x_{[c,d]})$ is a hypersurface, with $S' = \bK[x_i,\ldots,x_j]$ as in $\Cref{lem: minimal free resolution of R_Delta}$ below, and we treat it separately in each statement of this section. Both Hilbert series computed in this section rest on a single resolution, which we record in the next lemma.
\begin{lemma}
\label{lem: minimal free resolution of R_Delta}
    Suppose that neither of $A,B$ contain the other. 
    Let 
    \[
    S' = \bK[x_i,\ldots,x_j]\cong S/(x_1,\ldots,x_{i-1},x_{j+1},\ldots,x_n) \]
    and set $f = x_{A},g = x_{B}\in S',$ with degrees $\#A,\#B,$ respectively. Writing $f' = \lcm(f,g)/f$ and $g'= \lcm(f,g)/g,$ 
    \begin{align}
        0\longrightarrow S'(-min(\#A + \#B, m))
        \xlongrightarrow{
            \begin{psmallmatrix}
            f' \\ -g' 
            \end{psmallmatrix}
        } 
        S'(-\#A)\oplus S'(-\#B)
        \xlongrightarrow{
            \begin{psmallmatrix}
            f & g
            \end{psmallmatrix}
        } 
        S'
        \longrightarrow S'/(f,g)
        \longrightarrow 0
    \end{align}
    is a minimal free resolution of $S'/(f,g) = R_{\Delta_w}.$
\end{lemma}

\begin{proof}
    We begin by noting that since neither $A$ nor $B$ contains the other, then $\lcm(f,g) = x_{A\cup B}$ has degree $\min(\#A+\#B,m).$
    Now write $h = \gcd(f,g),$ so $f = hg'$ and $g = hf'$ with $\gcd(f',g') = 1.$
    Since $S'$ is a UFD, then if $\alpha f + \beta g = 0,$ then $\alpha g' = -\beta f'.$ 
    Furthermore, since $f',g'$ are coprime, then $\alpha = cf'$ and $\beta = -cg'$ for some $c\in S'.$
    Thus, the syzygy module of $(f,g)$ is free of rank one and generated by $(f',-g')$ and the sequence is exact at the middle term. Furthermore, since $S'$ is a domain, then we have injectivity at the left-hand map.
    The sequence is minimal: since $A\not\subseteq B$ and $B\not\subseteq A,$ then both $f'$ and $g'$ have positive degree, so every entry of both matrices lies in the graded maximal ideal. 
    Finally, if the two intervals overlap then neither is nested in the other, so together they reach endpoints $i,j$ and $A\cup B = [i,j],$ an interval of size $m.$ 
    However, if they are disjoint then $\#(A\cup B) = \#A + \#B.$ In both cases, $\#(A\cup B) = \min(\#A + \#B, m).$
\end{proof}


\begin{theorem}
\label{thm: hilbert series - single Schubert variety}
Let $w \in \sn(\calP)$, and let $m, \#A,$ and $\#B$ be as previously defined. Then the Stanley-Reisner ring $R_{\Delta_w} = S/I_{\Delta_w}$ has Hilbert series
\[
  \hilb\big(R_{\Delta_w},\, t\big)
  = \frac{1 - t^{\#A} - t^{\#B} + t^{\min(\#A + \#B,\, m)}}{(1-t)^{m}}
\]
Equivalently:
\begin{itemize}
\item if the interval nonfaces are disjoint \textup{(}$\#A + \#B \leq m$, i.e.\ $p < q$\textup{)},
\[
  \hilb\big(R_{\Delta_w},t\big) = \frac{(1-t^{\#A})(1-t^{\#B})}{(1-t)^{m}};
\]
\item if they overlap \textup{(}$\#A + \#B > m$, i.e.\ $p \leq q$\textup{)},
\[
  \hilb\big(R_{\Delta_w},t\big) = \frac{1 - t^{\#A} - t^{\#B} + t^{m}}{(1-t)^{m}};
\]
\item in particular, in the nested case ($\#A = m$ or $\#B = m$), with $k = \min(\#A,\#B),$
\[
\hilb\big(R_{\Delta_w},t) = \frac{1 - t^k}{(1-t)^m}.
\]
\end{itemize}
For $w = \mathrm{id}$ we have $\Delta_w = \{\emptyset\}$, so that
$I_{\Delta_w} = (x_1,\dots,x_n)$, $R_{\Delta_w} = \bK$, and
$\hilb(R_{\Delta_w},t) = 1$.
\end{theorem}

\begin{proof}
    In the nested case $R_{\Delta_w} = S'/(x_{[c,d]})$ is a hypersurface of degree $k$ in $m$ variables, so $\hilb\big(R_{\Delta_w},t) = (1 - t^k)/(1-t)^m,$ since either $\#A = m $ or $\#B = m.$
    Now suppose neither interval contains the other.
    Taking the alternating sum of the Hilbert series along the minimal free resolution described in \Cref{lem: minimal free resolution of R_Delta} and using $\hilb(S',t) = (1-t)^{-m}$ gives the displayed formula. In the disjoint case $\min(\#A + \#B,m) = \#A + \#B$ and $1 - t^{\#A} - t^{\#B} + t^{\#A + \#B} = (1-t^{\#A})(1-t^{\#B}),$ which gives us (a). 
    We note that this would also be clear without appealing to the minimal free resolution. 
    Since $f,g$ share no variables, then $g$ is a nonzerodivizor on $S'/(f),$ so $(f,g)$ forms a regular sequence, meaning $S'/(f,g)$ is a complete intersection, giving us the result.
    Now, in the overlapping case, $\min(\#A + \#B, m) = m,$ giving (b). The case for $w = \id$ is immediate.
\end{proof}

The resolution in \Cref{lem: minimal free resolution of R_Delta} also allows us to compute the depth and dimension of $R_{\Delta_w}.$
\begin{corollary}
\label{cor: len-depth-dim R_Delta}
    Let $w\neq\id\in\sn(\calP).$ The minimal free resolution of $R_{\Delta_w}$ over $S$ has 
    \begin{itemize}
        \item length $n - m +2$ (length $n - m + 1$ in the nested case)
        \item $\mathrm{depth}\,R_{\Delta_w} = m - 2$  ($\mathrm{depth}\,R_{\Delta_w} = m - 1$ in the nested case) 
        \item $\dim R_{\Delta_w} = 
        \begin{cases}
        m - 2 & p < q \\
        m - 1 & q\leq p.
        \end{cases}
        $
    \end{itemize}
    Consequently, when $p < q$ or in the nested case, $R_{\Delta_w}$ is Gorenstein. When $q\leq p$ and neither interval contains the other, then it is one short of being Cohen-Macaulay. 
\end{corollary}
\begin{proof}
    In the nested case, $R_{\Delta_w} = S'/(x_{[c,d]})$ and the minimal free resolution over $S$ is the Koszul complex on the regular sequence $x_1,\ldots,x_{i-1},x_{j+1},\ldots,x_n,x_{c,d]},$ of length $n-m+1;$ so $R_{\Delta_w}$ is a complete intersection, hence Cohen-Macaulay and Gorenstein, with $\depth R_{\Delta_w} = \dim R_{\Delta_w} = n - (n - m + 1) = m - 1.$ 
    Now suppose neither interval contains the other.
    The length is $n - m +2,$ since the minimal free resolution over $S,$ formed by tensoring the minimal free resolution from \Cref{lem: minimal free resolution of R_Delta} with a Koszul complex has $n-m$ generators $x_1,\ldots,x_{i-1},x_{j+1},\ldots,x_n.$ 
    The Auslander-Buchsbaum formula gives us that the depth is $n - (n - m + 2) = m - 2,$ and $\dim R_{\Delta_w} = \dim\Delta_w+1$ is given by \Cref{cor: purity of doubly pruned complexes}. 
    Finally, it is Gorenstein when $p < q$ since it is both Cohen-Macaulay ($\mathrm{depth}R_{\Delta_w} = \dim R_{\Delta_w}$) and a complete intersection ring. 
    When $q\leq p,$ $\mathrm{depth}R_{\Delta_w} = \dim R_{\Delta_w}-1.$
\end{proof}

As of this point, we have computed everything for the Stanley-Reisner ring of $\Delta_{w},$ $w\in\sn(\calP).$ We mostly do this for completeness. However, the ring that is more relevant to this paper is $R_{\Delta_w,n} = S/(I_\Delta+I_n).$ These are not simply a choice forced by the Borel presentation of the cohomology ring $H^*(G/B),$ but instead are a natural choice to reduce any Stanley-Reisner ring. Consider the image of the elementary symmetric functions generating $I_n$ under $S/I_{\Delta_w}$ -- these are exactly the \emph{universal system of parameters}. This system of parameters have made appearances in throughout the literature, but were initially named by Herzog and Moradi, and simultaneously defined by Adams and Reiner.
\begin{definition}[\cite{herzog-moradi,adams-reiner}]
\label{def: universal system of parameters}
    For a simplicial complex $\Delta$ of dimension $d-1,$ the \emph{universal system of parameters} $\theta = (\theta_1,\ldots,\theta_d)$ of $R_\Delta$ are 
    \[
    \theta_k = \sum_{\substack{F\in\Delta \\ \#F = k}}x_F\in R_\Delta,\quad k = 1,\ldots,d.
    \]
    They form a homogeneous system of parameters.
\end{definition}
Since $R_{\Delta_w,n} = S/(I_\Delta+I_n) = R_{\Delta_w}/(\theta),$ we are able to use this presentation to compute its Hilbert series. Since $R_{\Delta_w,n}$ is isomorphic to cohomology ring of the Schubert variety $X_w$ and \Cref{lem:defined by inclusions}, we translate the Hilbert series of this ring into a sum over the Bruhat interval below $w:$ since by \Cref{prop: schuberts form a basis for Rn} and since $J_w$ is spanned by Schubert polynomials \Cref{eq: Jw} the ring $R_n/J_w$ has a $\bK$-basis the images of the Schubert polynomials $\fraks_v$ with $v\leq w,$ and $\deg\fraks_v = \ell(v),$ we have
\[
\hilb(R_{\Delta_w,n},t) = \sum_{v\leq w}t^{\ell(v)}.
\]
Note that $\hilb(\bK[\theta_1,\ldots,\theta_d],t) = \prod_{r=1}^d(1-t^r)^{-1}.$
\begin{theorem}
\label{thm: hilbert series for R_Delta_n}
    Let $w\in\sn(\calP),$ $w\neq\id,$ with $m, \#A, \#B$ as above, with overlap $\gamma = \#A + \#B - m.$
    \begin{enumerate}
        \item If $p < q$ (disjoint, $\gamma\leq 0$), then $\theta_1,\ldots,\theta_{m-2}$ is a regular sequence on $R_{\Delta_{w}},$ and 
        \[
        \hilb(R_{\Delta_{w,n}},t) = [\#A]_t[\#B]_t[m-2]_t!.
        \]
        \item If $q\leq p$ and neither interval contains the other (overlapping, $i < q\leq p < j,$ $\gamma\geq 1$), then $\theta_1,\ldots,\theta_{m-1}$ is not a regular sequence on $R_{\Delta_w}$ and 
        \[
        \hilb(R_{\Delta_{w,n}},t) 
        = [\gamma]_t[m-1]_t! + t^\gamma~[m-\#B]_t[m-\#A]_t[m-2]_t!.
        \]
        \item If $q = i$ or $p = j$ (nested, $\#A = m$ or $\#B = m$), then $\theta_1,\ldots,\theta_{m-1}$ is a regular sequence on $R_{\Delta_w}$ and, with $ k = \min(\#A,\#B),$
        \[
        \hilb(R_{\Delta_{w,n}},t) = [k]_t[m - 1]_t!. 
        \]
    \end{enumerate}
\end{theorem}
\begin{proof}
    \begin{enumerate}
        \item By \Cref{cor: len-depth-dim R_Delta}, $R_{\Delta_w}$ is Cohen-Macaulay of Krull dimension $m - 2.$ 
        Thus it's a homogeneous system of parameters $\theta_1,\ldots,\theta_{m-2}$ is a regular sequence and 
        \begin{align*}
            \hilb(R_{\Delta_{w,n}},t) 
                & = \hilb(R_{\Delta_{w}},t)\prod_{r=1}^{m-2}(1-t^r) \\
                & = \frac{(1-t^{\#A})(1-t^{\#B})}{(1-t)^m}\cdot(1-t)^{m-2}[m-2]_t! \\
                & = [\#A]_t[\#B]_t[m-2]_t!.
        \end{align*}
        In the string of equalities, we use the standard equalities
        \[
        1-t^p = (1-t)[\#A]_t,
        \quad
        1-t^q = (1-t)[\#B]_t,
        \quad
        \prod_{r=1}^{m-2}(1-t^r) = (1-t)^{m-2}[m-2]_t!.
        \]
        
        \item In this case, $\dim R_{\Delta_w} = m - 1$ while $\depth\, R_{\Delta_w} = m - 2$ by \Cref{cor: len-depth-dim R_Delta}, so there is no homogeneous system of parameters of length $m - 1$ that also forms a regular sequence. 
        For the Hilbert series, recall that in the overlapping case $i<q\leq p<j,$ since neither nonface contains the other. Our strategy is as follows. We find an ideal $Z\subseteq R_{\Delta_w}$ such that both $Z$ and the quotient $R_{\Delta_w}/Z$ are, up to a shift, face rings of complexes whose Hilbert series modulo $\theta$ we already know. We also require of $Z$ that $R_{\Delta_w}/Z$ is Cohen-Macaulay, hence a free $\bK[\theta]$-module. The short exact sequence $0\to Z\to R_{\Delta_w}\to R_{\Delta_w}/Z\to 0$ then remains exact after passing modulo $\theta,$ and $\hilb(R_{\Delta_w,n},t)$ is the sum of the Hilbert series $Z/\theta Z$ and of $(R_{\Delta_w}/Z)/\theta(R_{\Delta_w}/Z).$
        Write 
        \[
        C := A\cap B=[q,p],
        \quad\#C = \gamma,
        \]
        and cover $\Delta_w$ by two subcomplexes, the faces that do not contain $C$ and the star of $C:$
        \[
        \Gamma_C := \{G\in\Delta_w:C\not\subseteq G\},
        \quad
        \star(C) := \{G:G\cup C\in\Delta_w\}
        \]
        so that $\Delta_w = \Gamma_C\cup\star(C)$ (the two subcomplexes are not disjoint; they meet in the faces $G\in\star(C)$ with $C\not\subseteq G).$ 
        Let $Z := (x_C)\subseteq R_{\Delta_w}$ be the ideal generated by the span of the monomials whose support contains $C,$ so that
        \[
        0 \longrightarrow 
        Z \longrightarrow 
        R_{\Delta_w}\longrightarrow 
        R_{\Delta_w}/Z \longrightarrow 
        0.
        \]
        First, we observe that both the quotient $R_{\Delta_w}/Z$ and the ideal $Z$ correspond to complexes that we have already seen in this paper.
        Beginning with the quotient, we notice that a monomial dies when it is divisible by $x_C,$ that is, when its support contains $C,$ thus $R_{\Delta_w}/Z = R_{\Gamma_C}.$     
        Since $C$ is contained in both $A$ and $B$, a subset of $[i,j]$ not containing $C$ contains neither $A$ nor $B$ and hence is a face $\Delta_w,$ so $\Gamma_C = \{G\subseteq[i,j]:C\not\subseteq G\}.
        $ and the minimal nonfaces of $\Gamma_C$ are $C$ together with the singletons outside $[i,j],$ with one minimal nonface of of cardinality of at least two.
        By \Cref{lem: two-min faces implies vertex-decomposable}, $\Gamma_C$ is vertex-decomposable, hence shellable. Moreover, by \Cref{prop: facet description of doubly pruned complexes} its facets are $[i,j]\setminus\{v\}$ with $v\in C,$ so it is pure of dimension $m - 2.$
        A pure shellable complex is Cohen-Macaulay, so $\theta_1,\ldots,\theta_{m-1}$ is a regular sequence on $R_{\Gamma_C}.$ 
        Thus, $R_{\Gamma_C} = S'/(x_C)$ is a hypersurface, where $S' = \bK[x_i,\ldots,x_j]$ as in \Cref{lem: minimal free resolution of R_Delta}, so $\hilb(R_{\Gamma_C},t) = (1-t^\gamma)/(1-t)^m.$ Since quotienting by a regular sequence of elements in degrees $1,\ldots, m - 1$ multiplies the Hilbert series by $\prod_{r = 1}^{m-1}(1-t^r),$ we obtain 
        \[
        \hilb(R_{\Gamma_C}/\theta R_{\Gamma_C},t) = 
        \frac{1-t^\gamma}{(1-t)^m}
        \prod_{r = 1}^{m - 1}(1-t^r)
        = [\gamma]_t[m - 1]_t!.
        \]

        We now move onto identifying the ideal $Z = (x_C).$
        The monomials of $Z$ are the $x_C x^{\vec{a}}$ with $C\cup\supp(\vec{a})\in\Delta_w,$ that is, with $\supp(\vec{a})\in\star(C).$ 
        Thus, division by $x_C$ is an isomorphism $Z\cong R_{\star(C)}(-\gamma)$ of graded vector spaces. 
        Moreover, it is also an isomorphism of modules. 
        For a face $F$ of $\Delta_w$ with $\#F = k$ and $g = x_C x^{\vec{a}}\in Z,$ the product $x_Fg$ survives exactly when $F\cup C\cup\supp(\vec{a})\in\Delta_w,$ so the terms of $\theta_k$ with $F\not\in\star(C)$ never survive. 
        Furthermore, for $F\in\star(C),$ the condition is simply the multiplication rule of $R_{\star(C)}.$ 
        Thus, for $G\subseteq[i,j]$ we have $G\cup C\in\Delta_w$ if and only if $[i,q-1]\not\subseteq G$ and $[p+1,j]\not\subseteq G.$ 
        That is, $\star(C)$ is the doubly pruned complex with support $[i,j]$ whose interval nonfaces are $[i,q-1]$ and $[p+1,j]$ of sizes $m - \#B$ and $m - \#A.$
        These are disjoint, since $q - 1 < p + 1,$ so the first case outlined in this theorem and already proved above can be applied to $\star(C).$ 
        This gives that
        \[\hilb(R_{\star(C),n},t) = [m-\#B]_t[m-\#A]_t[m-2]_t!.\]
        In particular, $\star(C)$ has Krull dimension $m - 2,$ so it has no faces of size $m - 1,$ and therefore $\theta_{m-1}$ acts on $Z$ as the empty sum $\theta_{m-1}(\star(C)) = 0.$ 
        Hence, $(\theta_1,\ldots,\theta_{m-1})Z = (\theta_1,\ldots,\theta_{m-2})Z,$ where the second sequence is the universal system of parameters of $R_{\star(C)}.$ 
        Therefore, $Z/\theta Z\cong (R_{\star(C)}/(\theta_1,\ldots,\theta_{m-2}))(-\gamma) = R_{\star(C),n}(-\gamma),$ and so
        \[
        \hilb(Z/\theta Z,t) = t^\gamma~[m-\#B]_t[m-\#A]_t[m-2]_t!.
        \]

        Finally, remembering that $R_{\Gamma_C}$ is Cohen-Macaulay and so a free $\bK[\theta]$-module, we 
        combine these Hilbert series calculations and identifications in order to reduce the earlier short exact sequence to another short exact sequence by considering it modulo $\theta:$
        \[
        0 \longrightarrow
        Z/\theta Z \longrightarrow
        R_{\Delta_w,n} \longrightarrow
        R_{\Gamma_C}/\theta R_{\Gamma_C} \longrightarrow
        0.
        \]
        Since Hilbert series are additive along short exact sequences, we obtain the required formula.

        \item By \Cref{cor: len-depth-dim R_Delta}, $R_{\Delta_w} = S'/(x_{[c,d]})$ is Cohen-Macaulay of Krull dimension $m - 1,$ and $\theta_r = 0$ for $r\geq m$ since $\Delta_w$ has no face of dimension $m - 1.$ Thus, its homogeneous system of parameters $\theta_1,\ldots,\theta_{m-1}$ is a regular sequence, and by \Cref{thm: hilbert series - single Schubert variety} 
        \[
        \hilb(R_{\Delta_{w,n}},t) = \hilb(R_{\Delta_w},t)\prod_{r=1}^{m-1}(1-t^r) = \frac{1-t^k}{(1-t)^m}\cdot(1-t)^{m-1}[m-1]_t! = [k]_t[m-1]_t!.
        \]
    \end{enumerate}
\end{proof}

Note that the formula of Part (2) of \Cref{thm: hilbert series for R_Delta_n} is valid in the nested case as well since there is one of $[m-\#A]_t,[m-\#B]_t$ is $[0]_t = 0$ and $\gamma = k.$ However, the argument of Part (2), is not, which is why we treat the proof for the nested case separately. For instance, when $w = s_r$ we have $m = k = 2$ and $\hilb(R_{\Delta_{w,n}},t) = [2]_t[1]_t! = 1 + t,$ the Poincar{\'e} polynomial of $X_{s_r} \cong \mathbb{P}^1,$ while for $w = w_0,$ the complex $\Delta_w$ is the boundary of the simplex on $[n], k = m = n,$ and $[n]_t[n-1]_t! = [n]_t!.$

\begin{example}
    Let $w = 45312\in\mathrm{S}_5.$ Then $w(3) = 3$ is the only fixed point. Then $i = 1, j = 5, p = w^{-1}(1) = 4,$ and $q = w^{-1}(5) = 2,$ so $A = [1,4]$ and $B = [2,5].$ 
    Since we are in the overlapping case: $q\leq p,$ then $ m = 5,\#A = \#B = 4,$ and $o = \#A + \#B - m = 3.$ \Cref{thm: hilbert series for R_Delta_n} gives
    \begin{align*}
        \hilb(R_{\Delta_{w,n}},t) 
        & = [3]_t[4]_t! + t^3[1]_t[1]_t[3]_t! \\
        & = 1 + 4t + 9t^2 + 15t^3 + 18t^4 + 16t^5 + 10t^6 + 4t^7 + t^8.
    \end{align*}
\end{example}

%------------------------------------------------------------
\subsection{The Adams-Reiner conjecture for doubly pruned simplical complexes}
%------------------------------------------------------------

Utilizing the universal system of parameters (\Cref{def: universal system of parameters}) Adams and Reiner posed a conjecture \cite[Conj.~6.1]{adams-reiner} that has only been proven for small families of simplicial complexes, but is open in general. It says that the shape of the minimal free resolution of $R_\Delta$ over the parameter ring $\bK[\theta]$ is governed by the barycentric subdivision $\Sd\Delta,$ colored so that the barycenter of a face $F$ of a simplicial complex $\Delta$ of dimension $d$ has color $\# F$ in $\Sd\Delta.$ Through their ``colorful" Hochster formula \cite[Thm.~3.3]{adams-reiner} and under their grading specialization $\epsilon_k\mapsto k,$ where $\epsilon_k$ is the $k$-th standard basis vector in $\bN^d$ it reads
\begin{equation}
\label{eq: Adams-Reiner conj}
    \dim_\bK
    \Tor_r^{\bK[\theta]}(R_\Delta,\bK)_{N} =
    \sum_{\substack{
    S\subseteq [d] \\
    \wt S = N
    }}
    \dim_\bK \Tilde{H}^{\# S - r - 1}(\Sd_S(\Delta);\bK).
\end{equation}
Here $d = \dim\Delta + 1,$ the weight $\wt S: = \sum_{k\in S}k$ records the internal degree, and $\Sd_S\Delta$ is the \emph{rank-selected} subcomplex of the barycentric subdivision $\Sd\Delta,$ (i.e., the order complex of $\{F\in\Delta:F\neq\varnothing,\# F\in S\}$).
We prove it for the first class where $R_\Delta$ fails to be free over $\bK[\theta]$ of simplicial complexes, namely those with at most two minimal nonfaces of cardinality of at least two, which by definition includes every doubly pruned complex, in particular, every $\Delta_w$. 
All (co)homology is reduced and taken with coefficients in $\bK.$ Since the complexes involved are finite and $\bK$ is a field, $\Tilde{H}^i(\delta;\bK)$ is dual to $\Tilde{H}_i(\Delta;\bK)$ by the universal coefficient theorem, so the two have the same dimension. Since \Cref{eq: Adams-Reiner conj} is only concerned with dimensions, we compute its RHS in homology, where the Mayer-Vietoris sequence is available.


We use only use three known facts within this section. First, \Cref{eq: Adams-Reiner conj} holds when $R_\Delta$ is Cohen-Macaulay \cite[Cor.~6.7]{adams-reiner} since both sides vanish for $r\geq 1$ giving for every Cohen-Macaulay complex $\Delta$ with $d = \dim\Delta + 1,$
\begin{equation}
\label{eq: conjecture in the CM case}
    \sum_{T\subseteq[d]}t^{\wt T}\dim\Tilde{H}_{\#T-1}(\Sd_T(\Delta); \mathbb{K}) = \hilb(R_{\Delta,n},t).
\end{equation}

Second, for such a Cohen-Macaulay complex $\Delta$ and $T\subseteq[d]$ the complex $\Sd_T(\Delta)$ is again Cohen-Macaulay of dimension $\# T - 1$ \cite{munkres,stanley-balanced}. In particular, $\Tilde{H}_i(\Sd_T(\Delta)) = 0$ for $i\neq\# T - 1.$ Third, the two sides of \Cref{eq: Adams-Reiner conj} have the same alternating sum over $r,$ both equal to $\hilb(R_\Delta,t)\prod_{r=1}^d(1-t^r)$ \cite[Prop~2.5,~Eq.~30]{adams-reiner}.

As stated before, a complex is a simplex, then its face ring is simply a polynomial ring. If the complex has exactly one minimal nonface, then its face ring is a hypersurface. If the complex has two disjoint minimal nonfaces, then its face ring is a complete intersection ring, and therefore Cohen-Macaulay. Thus, for the rest of this subsection, we can assume that the complex has two minimal nonfaces, each of cardinality of at least two, that are not disjoint (so they are overlapping). Call these sets $A$ and $B$ and call the overlap $C := A\cap B\neq\varnothing.$ Write $W = \supp(\Delta), m = \#W, \gamma = \#C, \alpha = \#(A\setminus B),$ and $\beta = \#(B\setminus A),$ all positive. Since $W\setminus\{v\}$ is a face exactly when $v\in C,$ then we have $\dim\Delta = m - 2$ and $d = m - 1.$ \Cref{rmk: conj into Delta w notation} translates this back into $\Delta_w$ notation.

The proof of \Cref{thm: hilbert series for R_Delta_n}(2) applies verbatim in this generality since it uses only that $A$ and $B$ are incomparable minimal nonfaces meeting in $C.$ 
The complexes $\Gamma_C,$ $\star(C),$ and the ideal $Z = (x_C)$ are defined there, and $\Delta = \Gamma_C\cup\star(C).$


Within $W = \supp(\Delta),$ the only minimal nonface of $\Gamma_C$ is $C.$
The complex $\star(C)$ has the two minimal nonfaces $A\setminus B$ and $B\setminus A;$ since $I_{\Gamma_{C}\cap\star(C)} = I_{\Gamma_C}+I_{\star(C)}.$ Finally, the complex $\Gamma_C\cap\star(C)$ has all three. 
Within each of these three lists, the sets are pairwise disjoint, so their monomials form a regular sequence and $R_{\Gamma_C}, R_{\star(C)},$ and $R_{\Gamma_C\cap\star(C)}$ are complete intersections, Cohen-Macaulay of Krull dimensions $d,d-1,$ and $d - 2,$ with
\begin{equation}
\label{eq: hilb of star and C intersect star}
    \hilb(R_{\star(C),n}, t) = [\alpha]_t[\beta]_t[m-2]_t!,
    \quad
    \hilb(R_{\Gamma_C\cap\star(C),n}, t) = [\gamma]_t[\alpha]_t[\beta]_t[m-3]_t!
\end{equation}

\begin{proposition}
\label{prop: dimensions formulas for reduced homologies}
    Let $S\subseteq [d]$ be nonempty, set $s = \#S$ and $S' = S\setminus\{d\},$ and $T = S\cap[d-2].$ Then $\Tilde{H}_i(\Sd_S(\Delta)) = 0$ for $i < s-2,$ and also for $i = s - 2$ when $d\not\in S.$ 
    \begin{itemize}
        \item If $d,d-1\in S,$ then
        \begin{align*}
        \dim\Tilde{H}_{s-2}\left(\Sd_S(\Delta)\right) 
        & = \dim\Tilde{H}_{s-2}\left(\Sd_{S'}(\star(C))\right) \\
        & + \dim\Tilde{H}_{s-3}\left(\Sd_{S'\setminus\{d-1\}}(\Gamma_C\cap\star(C))\right), 
    \end{align*}
    \item If $d\in S$ and $d - 1\not\in S,$ then
    \begin{align*}
        \dim\Tilde{H}_{s-2}\left(\Sd_S(\Delta)\right) 
        & = \dim\Tilde{H}_{s-2}\left(\Sd_{S'}(\star(C))\right) \\
        & - \dim\Tilde{H}_{s-2}\left(\Sd_{S'}(\Gamma_C\cap\star(C))\right).
    \end{align*}
    \end{itemize}
\end{proposition}
\begin{proof}
    A chain of faces of $\Delta$ with largest face $M$ lies in $\Gamma_C$ when $C\not\subseteq M$ and in $\star(C)$ when $C\subseteq M.$ 
    Thus, $\Sd_S(\Gamma_C)\cup\Sd_S(\star(C))$ with intersection $\Sd_S(\Gamma_C\cap\star(C)),$ and so we can apply Mayer-Vietoris.
    Moreover, $\star(C)$ has no face of size $d$ and $\Gamma_C\cap\star(C)$ has no face of size $d - 1,$ so $\Sd_S(\star(C)) = \Sd_{S\cap[d-1]}(\star(C))$ and $\Sd_S(\Gamma_C\cap\star(C)) = \Sd_T(\Gamma_C\cap\star(C)).$

    By the second fact listed a priori, the homology of the three complexes is thus concentrated in the single degrees $s-1,$ $\#(S\cap[d-1])-1,$ and $\#T-1,$  each between $s - 3$ and $s - 1.$
    For $i < s - 2,$ both groups on either side of $\Tilde{H}_i(\Sd_S(\Delta))$ vanish, which proves our first claim.
    If $d\not\in S,$ the degrees are $s - 1,$ $s - 1,$ and $s - 2$ or $s - 1.$ So, the groups on either side of $\Tilde{H}_{s - 2}(\Sd_S(\Delta))$ vanish as well.

    Let $d\in S$ so that $\# S' - 1 = s - 2.$
    \begin{itemize}
        \item If $d - 1\in S,$ then the three degrees are $s - 1, s - 2,$ and $s - 3$ so the sequence in degree $s - 2$ is
        \[
        0 
        \longrightarrow
        \Tilde{H}_{s-2}\left(\Sd_{S'}(\star(C))\right)
        \longrightarrow
        \Tilde{H}_{s-2}\left(\Sd_S(\Delta)\right)
        \longrightarrow
        \Tilde{H}_{s-3}\left(\Sd_T(\Gamma_C\cap\star(C))\right)
        \longrightarrow
        0.
        \]
        \item If $d-1\not\in S,$ then $S\cap[d-1] = T = S',$ and the degrees are $s-1,s-2,$ and $s-2,$ so the sequence in degree $s-2$ is
        \[
        \Tilde{H}_{s-2}\left(\Sd_T(\Gamma_C\cap\star(C))\right)
        \xlongrightarrow{\phi}
        \Tilde{H}_{s-2}\left(\Sd_{S'}(\star(C))\right)
        \longrightarrow
        \Tilde{H}_{s-2}\left(\Sd_S(\Delta)\right)
        \longrightarrow
        0,
        \]
        with $\phi$ induced by inclusion. 
        Every simplex of $\Sd_{S'}(\star(C))$ is a chain with cardinalities in $S'$ and so has dimension of at most $s - 2.$
        Thus, every $(s-1)$-simplex of $\Sd_S(\Delta)$ already lies in $\Sd_S(\Gamma_C).$
        Since neither complex has a simplex of dimension $s,$ then $\Tilde{H}_{s-1}(\Sd_S(\Gamma_C))\to\Tilde{H}_{s-1}\left(\Sd_S(\Delta)\right)$ is surjective. The connecting map out of $\Tilde{H}_{s-1}(\Sd_S(\Delta))$ is therefore zero, which means $\phi$ is injective, and $\Tilde{H}_{s-2}(\Sd_S(\Delta)) = \mathrm{coker}\phi.$
    \end{itemize}
\end{proof}



\begin{theorem}
\label{thm: adams-reiner conj for doubly pruned complexes}
    The Adams-Reiner conjecture (restated in \Cref{eq: Adams-Reiner conj}) holds over any field for every simplicial complex with at most two minimal nonfaces of cardinality at least two. In particular, for every doubly pruned complex, and so for every $\Delta_w$ with $w\in\sn(\calP).$ Moreover, $\Tor_r^{\bK[\theta]}(R_\Delta,\bK) = 0$ for $r\geq 2,$ and $\Tor_1$ vanishes unless the two nonfaces meet. In the latter case,
    \[
    \sum_N
    \dim_\bK \Tor_1^{\bK[\theta]}(R_\Delta,\bK)_N t^N 
    = t^{m-1+\gamma}~[\alpha]_t[\beta]_t[m-2]_t!.
    \]
\end{theorem}
\begin{proof}
    Call $P_r$ and $Q_r$ the generating functions in $t^N$ of the LHS and the RHS of \Cref{eq: Adams-Reiner conj}, respectively. We have completed the Cohen-Macaulay cases above, so we progress by assuming the remaining hypothesis. In this proof, we compute the LHS and the RHS, we show all terms on both sides are $0$ for $r\geq 2,$ and then compare the remaining generating functions. 
    
    \textbf{LHS:} As $\Gamma_C$ is Cohen-Macaulay of Krull dimension $d,$ it is $\bK[\theta]$-free, so the long exact sequence induced on $0\rightarrow Z\rightarrow R_\Delta\rightarrow R_{\Gamma_C}\rightarrow 0$ gives that $\Tor_r^{\bK[\theta]}(R_\Delta,\bK)\cong\Tor_r^{\bK[\theta]}(Z,\bK)$ for $r\geq 1,$ where $Z = (x_C)\cong R_{\star(C)}(-\gamma)$ with $\theta'$ acting as the universal system of parameters of $\star(C)$ and $\theta_d$ as zero. 
    Since $R_{\star(C)}$ is $\bK[\theta']$-free, the Koszul complex on $(\theta',\theta_d)$ has acyclic rows and zero differential in the $\theta$ direction. Thus, $\Tor_{\geq 2}(Z,\bK) = 0$ and $\Tor_1(Z,\bK)\cong (R_{\star(C)}/(\theta'))(-\gamma-d).$
    Hence, $P_r = 0$ for $r\geq 2,$ and by \Cref{eq: hilb of star and C intersect star}
    \begin{equation}
        P_1(t) = t^{\gamma+d}~\hilb(R_{\star(C),n},t) = t^{m-1+\gamma}~[\alpha]_t[\beta]_t[m-2]_t!\neq 0.
    \end{equation}

    \textbf{RHS:} By \Cref{prop: dimensions formulas for reduced homologies}, $Q_r = 0$ for $r\geq 2.$ Furthermore, we know that the only those $S$ containing $d$ contribute to $Q_1.$ Write such an $S$ as $S'\sqcup\{d\}$ with $S'\subseteq[d-1]$ so that $t^{\wt S} = t^dt^{\wt S'}.$
    Summing the first term of $\Cref{prop: dimensions formulas for reduced homologies}$ over these $S$ gives a $t^d$ term times a sum over all $S'\subseteq [d-1].$ In the second term, those $S'$ that contain $d-1$ become those subsets $T=S'\setminus\{d-1\}$ of $[d-2],$ which carry a factor of $t^{d-1},$ with a positive sign, while those avoiding $d-1$ are already such a $T$ and carry a negative sign. Together, they contribute the factor $t^{d-1}-1.$ Hence,
    \begin{align*}
        Q_1(t) = 
        & t^d~\sum_{S'\subseteq[d-1]}t^{\wt S'}
        \dim\Tilde{H}_{\#S'-1}\left(\Sd_{S'}(\star(C))\right) \\
        & + t^d(t^{d-1}-1)\sum_{T\subseteq [d-2]}t^{\wt T}
        \dim\Tilde{H}_{\#T-1}\left(\Sd_T(\Gamma_C\cap\star(C))\right).
    \end{align*}
    Both sums are instances of \Cref{eq: conjecture in the CM case}, so by \Cref{eq: hilb of star and C intersect star} they can be computed explicitly. 
    Using $1 - t^{d-1} = (1-t)[m-2]_t,$ then $[m-2]_t[m-3]_t! = [m-2]_t!$ and $(1-t)[\gamma]_t = 1 - t^\gamma.$ So,
    \[
    Q_1(t) = t^d~[\alpha]_t[\beta]_t[m-2]_t!
    \left(1-(1-t)[\gamma]_t\right)
    = t^{d+\gamma}~[\alpha]_t[\beta]_t[m-2]_t!,
    \]
    which is precisely $P_1(t).$ 
    Thus, $P_1 = Q_1,$ and since $P_r = Q_r = 0$ for $r\geq 2$ the third fact from above forces $P_0 = Q_0,$ as well.
\end{proof}

\begin{remark}
\label{rmk: conj into Delta w notation}
Let $w\in\sn(\calP),$ $w\neq\id,$ give $\Delta_w$ that sits in the overlapping case where $i < q \leq p < j,$ i.e., $q\leq p$ and neither interval contains the other, so that $A$ and $B$ are the two minimal nonfaces of $\Delta_w$ of cardinality of at least two. 
Discard any singleton nonfaces, which do not affect either side of \Cref{eq: Adams-Reiner conj}. 
Then $\Delta_w$ satisfies the hypothesis with $A = [i,p], B = [q,j], C = [q,p], W = [i,j],$ and $m = j - i + 1.$ 
Since incomparability forces $i < q\leq p < j,$ then
\[
\alpha = \#[i,q-1] = m - \#B,
\quad
\gamma = \#[q,p],
\quad
\beta = \#[p+1,j] = m - \#A.
\]
Thus, \Cref{thm: adams-reiner conj for doubly pruned complexes} proves \Cref{eq: Adams-Reiner conj} for $\Delta_w.$ Its row $r = 0$ is Part (2) of \Cref{thm: hilbert series for R_Delta_n} can now be recovered from the rank-selected homology of $\Sd\Delta_w$ rather than from the ideal $(x_C).$ Its row $r=1$ is new and is
\[
\sum_N
\dim_\bK\Tor_1^{\bK[\theta]}(R_{\Delta_w},\bK)_N~t^N
=
t^{m-1+\gamma}~[m-\#B]_t[m-\#A]_t[m-2]_t!
\]
and has total dimension $(m-\#B)(m-\#A)(m-2)!.$ 
By our calculation of $P_1(t)$ in the previous proof, it measures the failure of $R_{\Delta_w}$ to be free over $\bK[\theta].$
\end{remark}